%-------------------------------------------------------------------------------
%	SECTION TITLE
%-------------------------------------------------------------------------------

\iflanguage{german}{
	\cvsection{Erfahrung}
}
{
	\cvsection{Experience}
}

%-------------------------------------------------------------------------------
%\taity AG
%-------------------------------------------------------------------------------
\cvsubsection{aity AG}

\begin{cvparagraph}
	\iflanguage{german}
	{
		aity AG ist ein Unternehmen in der Finanzbranche, das Schweizer Engineering für Banking, SaaS und Embedded Finance bietet.\newline
		Das Unternehmen ist nach Unfix Patterns organisiert, einem nicht-traditionellen Modell für die Zusammenarbeit von Teams.
	}
	{
		aity AG is a company in the finance industry, providing Swiss engineering for banking, SaaS, and embedded finance.\newline
		The company is organised around Unfix patterns, a non-traditional model for how teams collaborate.
	}
\end{cvparagraph}

\begin{cventries}
	\cventry
	{
		\iflanguage{german}
			{Team Coach und Organisationsentwickler (Technisches Agile Coaching)}
			{Team Coach \& Organisational Developer (Technical Agile Coaching)}
	} % Job title
	{aity AG} % Organization
	{\iflanguage{german}{Bern, Schweiz}{Bern, Switzerland}} % Location
	{\iflanguage{german}{Mai 2025 - heute}{May 2025 - Present}} % Date(s)
	{
		\begin{cvitems}
			\item {
				\iflanguage{german}
				{Pair- und Mob-Programming direkt mit den Entwicklern, gemeinsam am Code statt nur am Whiteboard.}
				{Pair and mob with engineers at the keyboard during technical coaching sessions, contributing code alongside the team.}
			}
			\item {
				\iflanguage{german}
				{Leite den teamübergreifenden Roll-out von Engineering-Praktiken wie TDD, CI/CD, Trunk-Based Development und KI-unterstütztem Entwickeln. Erarbeite mit den Teams Branching-Strategien und gemeinsame Engineering-Qualitätsstandards.}
				{Lead cross-team rollout of engineering practices including TDD, CI/CD, trunk-based development, and AI-augmented coding. Work with teams to define branching strategies and shared engineering quality standards.}
			}
			\item {
				\iflanguage{german}
				{Etabliere Flow-Metriken und datenbasierte Entscheidungsfindung in der Engineering-Organisation, mit Forecasting und kontinuierlicher Verbesserung als Teil des Team-Alltags.}
				{Establish flow metrics and data-informed decision making across the engineering organisation, bringing forecasting and continuous improvement into team rhythms.}
			}
			\item {
				\iflanguage{german}
				{Arbeite mit Dutzenden von Entwicklern über mehrere Teams hinweg, gemeinsam mit weiteren Coaches und der Geschäftsleitung.}
				{Work with dozens of engineers across multiple teams, together with fellow coaches and leadership.}
			}
		\end{cvitems}
	}
\end{cventries}

\vspace{2em}

%-------------------------------------------------------------------------------
%\tLetPeopleWork GmbH
%-------------------------------------------------------------------------------
% Keep the company intro together with its role entry across page breaks
\needspace{16\baselineskip}
\cvsubsection{LetPeopleWork GmbH}

\begin{cvparagraph}
	\iflanguage{german}
	{
		LetPeopleWork GmbH unterstützt Unternehmen dabei, ihr volles Potenzial zu entfalten. Durch Beratung, Trainings und die Entwicklung von \underline{\href{https://letpeople.work/lighthouse}{Lighthouse}}, einem Tool für Flow-Metriken und datenbasierte Prognosen.\newline
		Wir teilen unser Wissen in unseren Fachgebieten, sammeln Erfahrungen in der Unternehmensführung und bleiben durch die Weiterentwicklung von \underline{\href{https://letpeople.work/lighthouse}{Lighthouse}} stets am Puls der Zeit.
	}
	{
		LetPeopleWork GmbH helps organizations unlock their full potential. Through consulting, training, and the development of \underline{\href{https://letpeople.work/lighthouse}{Lighthouse}}, a tool for flow metrics and data-driven forecasting.\newline
		We share knowledge in our areas of expertise, gain experience running a business, and stay up to date with technology by driving \underline{\href{https://letpeople.work/lighthouse}{Lighthouse}} development.
	}
\end{cvparagraph}

\begin{cventries}
	\cventry
	{
		\iflanguage{german}
			{Mitgründer, Geschäftsführer und Lead Developer Lighthouse}
			{Co-Founder, Managing Director and Lead Developer Lighthouse}
	} % Job title
	{
		LetPeopleWork GmbH
	} % Organization
	{
		\iflanguage{german}
		{ Zürich, Schweiz }
		{ Zurich, Switzerland }		
	} % Location
	{
		\iflanguage{german}{Mai 2025 - heute}{May 2025 - Present}
	} % Date(s)
	{
		\begin{cvitems}
			\item {
				\iflanguage{german}
				{Alleinige Konzeption, Entwicklung und Betrieb von \underline{\href{https://letpeople.work/lighthouse}{Lighthouse}} (C\#/.NET-Backend, React-/TypeScript-Frontend), aufgebaut in rund 20\% meiner Arbeitszeit. Lighthouse wird von Teams weltweit eingesetzt.}
				{Sole designer, developer, and operator of \underline{\href{https://letpeople.work/lighthouse}{Lighthouse}} (C\#/.NET backend, React + TypeScript frontend), built and maintained in roughly 20\% of my working time. Used by teams worldwide.}
			}
			\item {
				\iflanguage{german}
				{Architektur als modularer Monolith nach Ports-and-Adapters (ASP.NET Core, EF Core, React-SPA, SignalR, REST-API mit Konnektoren zu Jira, Azure DevOps, Linear und ServiceNow). Eine einzige Codebasis läuft als eigenständige Desktop-App ohne externe Abhängigkeiten, als selbst gehosteter Docker-Server mit PostgreSQL und als SaaS auf unserem eigenen Kubernetes-Cluster (Helm, ArgoCD), bei Bedarf mit mehreren Replikas (Redis, PostgreSQL-Advisory-Locks). Bewusst pragmatisch auf die Zielgruppe zugeschnitten: von Fachleuten, die Docker nicht kennen, bis zu Organisationen, die Helm-Charts erwarten.}
				{Architected as a ports-and-adapters modular monolith (ASP.NET Core, EF Core, React SPA, SignalR, REST API with connectors to Jira, Azure DevOps, Linear, and ServiceNow). One codebase runs as a zero-dependency standalone desktop app, as a self-hosted Docker server with PostgreSQL, and as SaaS on our own Kubernetes cluster (Helm, ArgoCD), with multiple replicas when needed (Redis, PostgreSQL advisory locks). Deliberately pragmatic for its audience: from people who have never touched Docker to organisations that expect Helm charts.}
			}
			\item {
				\iflanguage{german}
				{Lighthouse ist der Ort, an dem ich praktiziere, was ich coache: strikte TDD-Praxis, vollständig automatisierte CI/CD-Pipelines und agentisches Software-Engineering mit nWave ermöglichen einen kontinuierlichen wöchentlichen Release-Rhythmus.}
				{Lighthouse is where I practise what I coach: strict TDD, fully automated CI/CD pipelines, and agentic software engineering with nWave sustain a continuous weekly release cadence.}
			}
			\item {
				\iflanguage{german}
				{Leite Strategie, Kundenprojekte und Beratung zu Flow-Metriken, Prognosen sowie agilen und Lean-Praktiken parallel zur Lighthouse-Entwicklung.}
				{Lead strategy, client engagements, and consulting on flow metrics, forecasting, and agile/lean practices alongside Lighthouse development.}
			}
		\end{cvitems}
	}
\end{cventries}

\vspace{2em}

\cvsubsection{ABB Power Grids / Hitachi Energy}

\begin{cvparagraph}
	\iflanguage{german}
	{
		Hitachi Energy (ehemals ABB Power Grids) ist ein weltweit führendes Technologieunternehmen, das sich der Förderung einer nachhaltigen Energiezukunft verschrieben hat. Das Unternehmen entwickelt innovative Lösungen für Stromnetze und Energiesysteme und ist auf Hochspannungstechnologien, digitale Netzautomatisierung und Energiespeicherung spezialisiert. Ziel ist es, die Umstellung auf erneuerbare Energien voranzutreiben und eine zuverlässige Stromversorgung weltweit sicherzustellen.
		
		Ich begann meine Laufbahn bei ABB während meiner Ausbildung und arbeitete während meines Software-Engineering-Studiums weiterhin auf Stundenbasis dort. Nach dem Abschluss meines Studiums trat ich Vollzeit als Software Engineer in das Unternehmen ein und habe im Laufe der Jahre verschiedene Rollen und Verantwortlichkeiten übernommen, die zu diversen Initiativen beigetragen haben.
	}
	{
		Hitachi Energy (formerly ABB Power Grids) is a global technology leader committed to advancing a sustainable energy future. The company develops innovative solutions for power grids and energy systems, specializing in high-voltage technologies, digital grid automation, and energy storage. Its mission is to drive the transition to renewable energy and ensure reliable power delivery worldwide.
		
		I began my journey at ABB during my apprenticeship, continuing to work on an hourly basis while pursuing my Software Engineering studies. After completing my degree, I joined the company full-time as a Software Engineer and have since taken on diverse roles and responsibilities, contributing to various initiatives over the years.
}
\end{cvparagraph}

%-------------------------------------------------------------------------------
%	CONTENT
%-------------------------------------------------------------------------------
\begin{cventries}
	\cventry
	{
		\iflanguage{german}
			{ Leiter der Prozesse und Qualität für Plattformen, Agile Coach (Senior Software Engineer II) }
			{ Head of Processes and Quality for Platforms, Agile Coach (Senior Software Engineer II) }
	} % Job title
	{Hitachi Energy} % Organization
	{
		\iflanguage{german}
		{ Baden, Schweiz }
		{ Baden, Switzerland }
	} % Location
	{
		\iflanguage{german}
			{Okt. 2023 - März 2025}
			{Oct. 2023 - Mar. 2025}
	} % Date(s)
	{
		\begin{cvitems}
			\item {
				\iflanguage{german}
				{
					Trieb Prozessverbesserungen über 10+ Teams hinweg voran, inklusive eingebundener externer Offshore-Teams. Erarbeitete gemeinsam mit Scrum Mastern und Product Ownern klare Team-Grenzen, die Selbstorganisation ermöglichten und gleichzeitig einen einheitlichen Arbeitsstil sowie gemeinsame Qualitätsstandards sicherstellten: klare Qualitätsziele, einen standardisierten Workflow, Definition-of-Done-Kriterien und zentrale Performance-Metriken.
				}
				{
					Drove process improvements across 10+ teams, including offshore externals integrated into the setup. Partnered with Scrum Masters and Product Owners to define clear team boundaries that empowered self-management while preserving a unified working style and shared quality standards: clear quality goals, a standardised workflow, Definition-of-Done criteria, and key performance metrics.
				}
			}			
			\item {
				\iflanguage{german}
				{
					Verantwortete die Qualität von drei internen Plattformen. Etablierte mit Entwicklern und Product Ownern objektive, datengetriebene Qualitätsmetriken anstelle subjektiver Bewertungen. Entwickelte Dashboards und automatisierte Datenerfassung für Team- und Stakeholder-Transparenz und unterstützte Teams dabei, aus diesen Daten validierte Verbesserungsexperimente abzuleiten.
				}
				{
					Owned quality for three internal platforms. Established objective, data-driven quality metrics with developers and Product Owners, replacing subjective assessment. Designed dashboards and automated data collection for team and stakeholder transparency, and supported teams in turning that data into validated improvement experiments.
				}
			}	
			\item {
				\iflanguage{german}
				{
					Hielt interaktive Scrum-Schulungen für Gruppen von bis zu 50 Personen aus verschiedenen Abteilungen, mit Fokus auf Verhaltensänderung statt Ritualen (Liberating Structures, Training from the Back of the Room).
				}
				{
					Delivered interactive Scrum training to groups of up to 50 people across departments, focused on behaviour change rather than ceremonies (Liberating Structures, Training from the Back of the Room).
				}
			}
			\item {
				\iflanguage{german}
				{
					Entwickelte eigene Tools zur Automatisierung von Projektmanagement- und Metrik-Aufgaben, darunter Portfolio-Prognosen für Features über mehrere Teams. Die probabilistischen Prognosen ersetzten aufwändige Schätzungen, waren genauer und wurden später zur Inspiration für Lighthouse.
				}
				{
					Built my own tools to automate project-management and metrics work, including portfolio-level forecasting for features spanning several teams. The probabilistic forecasts replaced effortful estimation, proved more accurate, and later inspired Lighthouse.
				}
			}
		\end{cvitems}
	}	
	
  	% Empty Entry to make sure there is enough whitespace between the sections
	\cventry
	{ 	}
	{  	}
	{  	}
	{  	}
	{  	}
	
	\cventry
	{
		\iflanguage{german}
			{ Tech Lead, Scrum Master, Configuration Manager (Senior Software Engineer) }
			{ Tech Lead, Scrum Master, Configuration Manager (Senior Software Engineer) }
	} % Job title
	{Hitachi Energy} % Organization
	{
		\iflanguage{german}
			{ Baden, Schweiz }
			{ Baden, Switzerland }
	} % Location
	{
		\iflanguage{german}
			{ Dez. 2019 - Okt. 2023 }
			{ Dec. 2019 - Oct. 2023 }
	} % Date(s)
	{
		\begin{cvitems}
			\item {
				\iflanguage{german}
				{
					Verantwortete Configuration Management und Pipelines für eine Produktlinie mit mehreren Produkten und Teams. Leitete die Modernisierung der CI-Umgebungen und automatisierte zuvor manuelle Prozesse wie digitales Signieren, Scanning auf Open-Source-Lizenzverletzungen und die Behebung von Cybersecurity-Schwachstellen. Gründete ein \glqq Configuration Management Forum\grqq{} zum Austausch bewährter Praktiken zwischen Teams und Produkten.
				}
				{
					Owned Configuration Management and pipelines for a Product Line spanning multiple products and teams. Led the modernisation of CI environments, automating previously manual processes (digital signing, open-source licence scanning, cybersecurity vulnerability remediation). Founded a `Configuration Management Forum' to share practices across teams and products.
				}
			}
			\item {
				\iflanguage{german}
				{
					Als Scrum Master leitete ich ein Team, das eine C\#/React/TypeScript-Anwendung auf der Azure Cloud entwickelte, und baute die vollständige CI/CD-Automatisierung selbst auf: vom mühsamen vierteljährlichen manuellen Release-Prozess (samt stressiger Hotfix-Schleifen) hin zu mehreren auslieferbaren Builds pro Tag und produktiven Releases auf Knopfdruck (im Schnitt monatlich). Wöchentliche \glqq Scrum Bites\grqq{} mit Product Owner und Entwicklern vertieften jeweils ein Konzept.
				}
				{
					As Scrum Master, led a team building a C\#/React/TypeScript application on Azure Cloud, and built the full CI/CD automation myself, shifting from a tedious quarterly manual release process (with stressful urgent-hotfix cycles) to multiple releasable builds per day and on-demand production deploys (roughly monthly in practice). Ran weekly `Scrum Bites' with the Product Owner and developers to evolve one practice at a time.
				}
			}
			\item {
				\iflanguage{german}
				{
					Als Scrum Master leitete ich ein Team durch die 8-monatige Modernisierung einer rund 350k LOC grossen Silverlight-Anwendung auf .NET 6+, TypeScript und React. Führte TDD und Clean-Code-Praktiken ein, sodass die Migration inkrementell ablief und für bestehende Nutzer risikoarm blieb.
				}
				{
					As Scrum Master, led a single team through the 8-month modernisation of a ~350k LOC Silverlight application onto .NET 6+, TypeScript, and React. Introduced TDD and clean-code practices so the rebuild stayed incremental and low-risk for existing users.
				}
			}
			\item {
				\iflanguage{german}
				{
					Coachte und trainierte mehrere Teams in Scrum und Kanban, führte Flow-Metriken und probabilistische Prognosen ein und gründete eine unternehmensweite \glqq Scrum Master Guild\grqq{}, die gute Praktiken über interne Veranstaltungen und Meetups verbreitete.
				}
				{
					Coached and trained multiple teams in Scrum and Kanban, introduced flow metrics and probabilistic forecasting, and founded a company-wide `Scrum Master Guild' that spread good practices through internal events and meetups.
				}
			}
		\end{cvitems}
	}	
	
  	% Empty Entry to make sure there is enough whitespace between the sections
	\cventry
	{ 	}
	{  	}
	{  	}
	{  	}
	{  	}
	
	\cventry
	{
		\iflanguage{german}
			{ Software Entwickler, Scrum Master (Software Engineer) }
			{ Software Developer, Scrum Master (Software Engineer) }
	} % Job title
	{ABB Power Grids} % Organization
	{
		\iflanguage{german}
			{ Baden, Schweiz }
			{ Baden, Switzerland }
	} % Location
	{
		\iflanguage{german}
			{ Aug. 2014 - Dez. 2019 }
			{ Aug. 2014 - Dec. 2019 }
	} % Date(s)
	{
		\begin{cvitems}
			\item {
				\iflanguage{german}
				{
					Entwickelte modulare Windows-Desktop-Anwendungen in C\#/.NET mit WPF/XAML.
				}
				{
					Developed modular Windows desktop applications in C\#/.NET with WPF/XAML.
				}
			}
			\item{
				\iflanguage{german}
				{
					Mitglied des \glqq Clean Code Teams\grqq{}, das die Einführung von statischer Codeanalyse, Pair Programming, Code-Reviews, Refactoring, Unit Testing und TDD vorantrieb. Hielt Einführungsschulungen und praxisorientierte Workshops, um Kollegen bei der täglichen Anwendung zu unterstützen.
				}
				{
					Member of the `Clean Code Team' driving adoption of static analysis, pair programming, code reviews, refactoring, unit testing, and TDD. Ran introductory training and hands-on workshops to help teammates embed these practices day-to-day.
				}
			}
			\item{
				\iflanguage{german}
				{
					Baute und erweiterte eine Continuous-Integration-Umgebung auf Azure Pipelines, mit Automatisierung von Builds, Testdurchläufen und Paketauslieferung in Testumgebungen für schnellere Feedback-Zyklen und zuverlässigere Releases.
				}
				{
					Built and extended a continuous integration environment on Azure Pipelines, automating builds, test runs, and package deployment to test environments for faster feedback and more reliable releases.
				}
			}
			\item{
				\iflanguage{german}
				{
					Baute ein automatisiertes End-to-End-Testframework auf Basis des Page-Object-Patterns, das gegen die aktuellsten Pipeline-Pakete läuft und Änderungen kontinuierlich validiert.
				}
				{
					Built an automated end-to-end test framework based on the page-object pattern, running against the latest pipeline packages to validate changes continuously.
				}
			}	
			\item{
				\iflanguage{german}
				{
					Leitete die Skalierung von 1 auf 5+ Teams in 5 Ländern. Führte mehrwöchige On-Site-Trainings und praktisches Onboarding durch und etablierte ein skaliertes Prozessmodell auf Basis von Nexus/SAFe.
				}
				{
					Led the scale-up from 1 to 5+ teams across 5 countries. Delivered several weeks of on-site training and hands-on onboarding, and established a Nexus/SAFe-based scaled process model.
				}
			}		
			\item {
				\iflanguage{german}
				{
					Begleitete einen Informatiklehrling durch die Ausbildung und unterstützte seine Entwicklung im Software Engineering im Tagesgeschäft.
				}
				{
					Mentored an application developer apprentice through their training, guiding their day-to-day growth in software engineering.
				}
			}  		
		\end{cvitems}
	}	
\end{cventries}
